% =============================================================================
% Ethics Statement and Broader Impact
% Self-contained file: % =============================================================================
% Ethics Statement and Broader Impact
% Self-contained file: % =============================================================================
% Ethics Statement and Broader Impact
% Self-contained file: % =============================================================================
% Ethics Statement and Broader Impact
% Self-contained file: \input{ethics_statement} into main.tex
% Intended placement: after the Conclusion section, before References.
% Label: \label{sec:impact_standalone}
% Target length: >= 2 pages (NeurIPS Broader Impact guidance).
% =============================================================================

\section{Ethics, Limitations, and Broader Impact}
\label{sec:impact_standalone}

This statement is written to satisfy the NeurIPS Code of Ethics and Broader
Impact requirements. It complements the shorter Limitations and Broader
Impact subsections embedded in Section~\ref{sec:limitations} by consolidating
in one place a full dual-use analysis, itemised compute accounting, carbon
footprint estimation, data provenance disclosures, a candid acknowledgment
of the closed-source training backend on which a fraction of our headline
numbers depends, and a list of methodological limits we are aware of but
did not have resources to close within the submission window.

% ---------------------------------------------------------------------------
\subsection{Dual-Use Analysis: Misuse Risks of Reasoning RL}
\label{sec:dualuse}
% ---------------------------------------------------------------------------

\paragraph{Threat model.}
\ours{} releases (i) training scripts that apply GRPO to the GSM8K
\cite{cobbe2021training} mathematical reasoning benchmark and to the
Salesforce xLAM-Function-Calling-60k \cite{liu2024xlam} tool-use corpus,
(ii) a set of LoRA adapters published on the HuggingFace Hub
(\texttt{arvindcr4/tinker-rl-bench-*}), and (iii) diagnostic code for
Zero-Variance Fraction, reward-stability, and length-bias analysis. We
analyse four concrete misuse pathways these artefacts could, in principle,
enable, and we describe the existing guardrails and the residual risk we
judge acceptable for publication.

\paragraph{M1. Weaponisation of mathematical reasoning.}
The most capability-relevant artefact we release is GRPO training code that
improves grade-school arithmetic and multi-step reasoning. A malicious
operator could attempt to extend this pipeline to tasks of greater
dual-use concern, e.g., chemistry olympiad problems, cryptographic key
recovery, or financial fraud. We judge the \emph{marginal} uplift provided
by our release to be small for three reasons. First, GRPO at our training
scale does not create new capabilities; the base-model control for GSM8K
(Section~\ref{sec:results}, 82.0\%) shows that the bulk of held-out
accuracy is attributable to Qwen3-8B's pretraining. Our full GRPO run only
adds \(+1.3\) percentage points ($p{=}0.26$, not statistically significant).
Second, the public literature already contains GRPO implementations
\cite{shao2024deepseekmath, openrlhf2024, verl2024, trl2020} that are at
least as capable. Third, any reasoning uplift is bounded by the rewarder:
GSM8K rewards use exact-match against a human-written gold answer, which
does not generalise to the domains listed above without a new reward
signal, which itself requires expertise and labelled data.

\paragraph{M2. Reward hacking and the long-term alignment risk.}
GRPO, like all policy-gradient RL from a proxy reward, is vulnerable to
reward hacking: the model learns to exploit idiosyncrasies of the reward
function rather than solve the intended task \cite{skalse2022specificationgaming,
krakovna2020specification}. Our length-bias analysis
(Section~\ref{sec:length_bias}) and the verbosity-trap finding in 4/11
GRPO runs are concrete demonstrations of this failure mode in our own data.
Users who apply our scripts to under-specified reward functions in
production settings (e.g., a custom ``helpfulness'' classifier) may observe
models that appear to improve on held-out metrics while silently
optimising against unintended proxies. We include the Zero-Variance
Fraction and reward-stability diagnostics precisely so that downstream
users can detect such drift.

\paragraph{M3. Circumventing safety training via distillation or tool-use.}
The tool-use track trains LoRA adapters that teach a base model to emit
schema-valid JSON function calls \cite{liu2024xlam}. A motivated adversary
could in principle use the same pipeline to teach a base model to emit
jailbreak prompts as ``tool calls'' or to invoke an adversarial external
tool. We mitigate this by (a) releasing only adapters, not merged weights,
so that the safety-trained instruct checkpoint must be applied separately
and any instruct-layer guardrails remain active; (b) documenting in model
cards that the tool-use adapters are trained on a narrow schema (five
synthetic tools) and have not been safety-evaluated for open-ended function
calling; and (c) not releasing any function-calling harness that would
actually execute emitted calls.

\paragraph{M4. Compute-efficiency externalities.}
Our work argues that critic-free RL with careful group-size selection
(\(G{=}32\)) can be run on modest budgets. Improved training efficiency is
itself dual-use: it lowers the cost of running adversarial fine-tuning at
scale. We judge this externality to be modest relative to the already-low
cost of fine-tuning a 0.6B--8B model with commodity LoRA recipes
\cite{hu2022lora}, but we flag it explicitly.

\paragraph{Residual risk.}
After weighing M1--M4 we judge that the reproducibility, calibration, and
pedagogical benefits of releasing \ours{} outweigh the residual misuse
risk. We adopt the standard mitigation of releasing adapters (not merged
weights), documenting intended use and out-of-scope use in model cards,
and publishing alongside the diagnostics a researcher would need to detect
reward hacking.

% ---------------------------------------------------------------------------
\subsection{Compute Cost Accounting}
\label{sec:compute_accounting}
% ---------------------------------------------------------------------------

Table~\ref{tab:compute_spend} itemises the financial cost of every
platform used in the project. Costs are real spend for the submitting
authors; no in-kind credits are omitted. Dollar figures are in USD.

\begin{table}[ht]
\centering
\small
\caption{Itemised compute spend across platforms. Totals include both
successful and failed runs, because failed-run spend is real and should
not be hidden from reviewers \cite{henderson2018deep}.}
\label{tab:compute_spend}
\begin{tabular}{llrr}
\toprule
\textbf{Platform} & \textbf{Use} & \textbf{Runs} & \textbf{Cost (USD)} \\
\midrule
Tinker SDK v0.16.1 & GRPO on 4B/8B (original suite)        & 25 & \$40--55 \\
Tinker SDK v0.16.1 & 10x Structural Ceiling ablation       & 32 & \$65 \\
Tinker SDK v0.16.1 & World-Class Suite (frontier/MoE)      & 13 & \$25--30 \\
Modal H100         & PPO baselines (Qwen3-8B, Llama-3.1-8B)& 2  & \$12--18 \\
Modal H100         & KL-tracking run (failed)              & 1  & \$2--4   \\
Google Colab Pro (T4) & 0.5B--3B QLoRA SFT+GRPO            & --- & \$10/person \\
NVIDIA L4 (TRL baseline)& Qwen2.5-0.5B GSM8K, 5 seeds      & 5  & \$3--5   \\
HuggingFace Hub    & Model + adapter hosting               & --- & \$0 \\
Weights \& Biases  & Experiment tracking (academic tier)   & --- & \$0 \\
\midrule
\multicolumn{3}{l}{\textbf{Total authors' out-of-pocket spend}} & \textbf{\$130--140} \\
\bottomrule
\end{tabular}
\end{table}

\paragraph{Hidden costs.} Beyond direct platform spend, the project
consumed human labour (approximately six student-FTE-months, unpaid) and
infrastructure generously subsidised by PES University's LLM Research
Group (shared access to one A100 node used for TRL baselines and
pre-submission sanity checks). No author received commercial compensation
from Thinking Machines, Modal, or any other vendor mentioned.

\paragraph{Failed-run accounting.} Of the Tinker spend, we estimate
\textasciitilde\$30--35 was consumed by runs that ultimately failed
(JWT expiry, API stalls for GPT-OSS-20b and Kimi-K2 before re-run,
KL-tracking gradient bug). We disclose this because reviewers often ask
whether reported total spend reflects only clean, successful runs; it does
not. Excluding failed spend would understate the true resource cost of
reproducing our work by roughly 25\%.

% ---------------------------------------------------------------------------
\subsection{Carbon Footprint Estimation}
\label{sec:carbon}
% ---------------------------------------------------------------------------

Table~\ref{tab:carbon2} computes an energy and CO$_2$-equivalent estimate
for each platform. We follow the methodology of
\citet{patterson2024empirical} and \citet{strubell2019energy}: for each
GPU-hour we multiply device TDP by wall-clock GPU-hours and the data-centre
Power Usage Effectiveness (PUE), then multiply by the grid carbon intensity
in the region where the hardware is believed to be located.

\paragraph{Assumed constants.}
\begin{itemize}
  \item NVIDIA H100 SXM5 TDP: 700\,W.
  \item NVIDIA A100 80GB TDP: 400\,W.
  \item NVIDIA L4 TDP: 72\,W.
  \item NVIDIA T4 TDP: 70\,W.
  \item Data-centre PUE: 1.1 (conservative; typical hyperscale PUE is
        1.10--1.15 \cite{google2023pue, aws2023pue}).
  \item US average grid intensity (EPA eGRID 2023): 0.367\,kg\,CO$_2$-eq / kWh
        \cite{epa2023egrid}.
  \item Indian average grid intensity (CEA 2023): 0.716\,kg\,CO$_2$-eq / kWh
        \cite{cea2023co2}.
  \item We use the US average for Modal H100 (US-East region as configured)
        and Tinker (location undisclosed; assumed US), and the Indian
        average for Colab Pro T4 used by authors based in Bengaluru
        (Google Cloud asia-south1).
\end{itemize}

\paragraph{GPU-hour estimates.}
Tinker GPU-hour counts are not disclosed by the platform. We estimate them
from wall-clock run duration and assume a single H100-class accelerator per
job, which is consistent with Tinker's published architecture for the
model sizes we trained. Modal GPU-hours are measured directly from Modal's
billing dashboard.

\begin{table}[ht]
\centering
\small
\caption{Energy and carbon footprint estimates. Tinker hardware is not
publicly disclosed; we assume 1$\times$ H100 equivalent per run. Failed /
stalled runs are included. Totals rounded to one significant figure.}
\label{tab:carbon2}
\begin{tabular}{lrrrrr}
\toprule
\textbf{Platform} & \textbf{GPU-h} & \textbf{TDP\,(W)} & \textbf{PUE} &
\textbf{Energy\,(kWh)} & \textbf{CO$_2$\,(kg)} \\
\midrule
Tinker (assumed H100)   & \textasciitilde 950 & 700 & 1.1 & \textasciitilde 732 & \textasciitilde 269 \\
Modal H100 (US-East)    & \textasciitilde 18  & 700 & 1.1 & \textasciitilde 14  & \textasciitilde 5.1 \\
PES A100 (in-kind)      & \textasciitilde 60  & 400 & 1.1 & \textasciitilde 26  & \textasciitilde 19   \\
Colab Pro T4 (asia-south1) & \textasciitilde 40 & 70 & 1.2 & \textasciitilde 3.4 & \textasciitilde 2.4  \\
NVIDIA L4 (TRL baseline)   & \textasciitilde 10 & 72 & 1.1 & \textasciitilde 0.8 & \textasciitilde 0.3  \\
\midrule
\textbf{Project total (best estimate)} & & & & \textasciitilde 776 & \textasciitilde 296 \\
\bottomrule
\end{tabular}
\end{table}

\paragraph{Interpretation and uncertainty.}
Our best estimate of \textasciitilde 296\,kg\,CO$_2$-eq for the entire
project is comparable to a single round-trip economy flight
Bengaluru--Delhi (\textasciitilde 300\,kg). The dominant uncertainty is
the Tinker GPU-hour count: because Tinker does not expose hardware telemetry,
a factor-of-two error in GPU-hours or TDP is plausible. Under a pessimistic
assumption (2$\times$ GPU-hours, H100 running near 100\% TDP continuously),
the Tinker contribution could be as high as \textasciitilde 540\,kg CO$_2$-eq.
Under an optimistic assumption (Tinker backends actually use more
energy-efficient accelerators than H100, 50\% average utilisation) the
contribution could be as low as \textasciitilde 130\,kg. We report the
central estimate in the main table and caution readers that it should not
be interpreted as precise. \emph{All} numbers should be regarded as upper
bounds on the carbon cost of reproducing our work, because subsequent
users can skip the failed runs and the ablation sweeps.

\paragraph{Offset and mitigation.}
We did not purchase carbon offsets. Instead, we have (a) released all
trained checkpoints on HuggingFace Hub so that downstream users do not
need to re-run our experiments; (b) released step-level CSV logs so that
learning curves can be inspected without re-training; and (c) documented
in Section~\ref{sec:compute_accounting} which runs were wasteful, so that
replicators can skip them. We regard reproducibility-by-artefact as a
more durable mitigation than offsets.

% ---------------------------------------------------------------------------
\subsection{Data Provenance}
\label{sec:provenance}
% ---------------------------------------------------------------------------

\ours{} uses only publicly released research datasets. No private data,
personally identifiable information, or licensed proprietary content is
used. No human annotators were employed during this work. We document
each dataset below.

\paragraph{GSM8K \cite{cobbe2021training}.}
8{,}500 grade-school math word problems (7{,}473 train / 1{,}319 test)
authored by human writers contracted by OpenAI. Released under the
\textbf{MIT License} via \url{https://github.com/openai/grade-school-math}.
We use the standard train/test splits without modification. The canonical
citation is \citet{cobbe2021training}. Known limitations of GSM8K include
gender-neutral but culturally US-centric word problems (names, currencies,
sports) and a moderate rate of human-labelling errors estimated at
\textasciitilde 2\% by \cite{zhang2024gsm1k}. Our held-out evaluation
respects the test/train split.

\paragraph{Salesforce xLAM-Function-Calling-60k \cite{liu2024xlam}.}
60{,}000 function-calling examples generated by Salesforce Research as
training data for the xLAM agentic model family. Released under the
\textbf{CC BY 4.0} license via \url{https://huggingface.co/datasets/Salesforce/xlam-function-calling-60k}.
We use the public release as-is; specifically, we use the first 35 prompts
$\times$ 10 rollouts in the 10x Structural Ceiling tool-use track and the
full 60k split for the xlam-60k real-data run in Section~\ref{sec:results}.
Known limitations: xLAM schemas are synthetic and skew toward cleanly-typed
arguments; the dataset under-represents ambiguous, error-handling, and
multi-turn tool-calling. Attribution to Salesforce is required by the
license and is provided in Section~\ref{sec:acknowledgments} and in the
xLAM-derived model cards.

\paragraph{HumanEval \cite{chen2021codex}.}
164 Python programming problems with unit tests, released by OpenAI under
the \textbf{MIT License}. Used unmodified for pass@$k$ evaluation. Known
limitations: small scale, English-language docstrings, single-language
coverage; documented contamination concerns against frontier pretraining
corpora \cite{riddell2024contamination}.

\paragraph{NuminaMath \cite{li2024numinamath}.}
Used by collaborator Mohammad Rafi for the multi-stage
GSM8K+NuminaMath pipeline. Released under \textbf{Apache 2.0}. We use a
downstream checkpoint reported by Rafi; our repository contains only his
scripts and the Apache-licensed derivative weights.

\paragraph{Open-Platypus \cite{lee2023platypus}.}
3{,}000 SFT examples used by collaborator Madhu for Qwen3-8B code
generation warm-up. Open-Platypus is a curated subset released under
\textbf{CC BY-NC 4.0}; the non-commercial restriction is respected in our
release, which is academic-use only.

\paragraph{Synthetic tool-use corpus (authored).}
We additionally generated a small five-tool synthetic corpus (calculator,
web-search stub, calendar, file-read, email-send) used for Section 4.1's
saturation analysis. The corpus is authored by the paper's authors from
publicly documented APIs, contains no third-party content, and is
released under the \textbf{MIT License} in our repository under
\texttt{data/synthetic\_tools/}. Generation prompts are included for
auditability. No user data, scraped web content, or proprietary API
traces are present.

\paragraph{No web scraping, no human subjects.}
We did not scrape any website. We did not run any study with human
participants. No IRB review was therefore required. No data that could
reasonably be considered to contain PII, copyrighted content, or sensitive
personal attributes was used at any stage of training, evaluation, or
reward modelling.

% ---------------------------------------------------------------------------
\subsection{Closed-Source Tinker Acknowledgment}
\label{sec:tinker_closed}
% ---------------------------------------------------------------------------

A central tension in \ours{} is that a substantial fraction of our
headline numbers come from a closed-source commercial platform while our
paper simultaneously advocates for reproducibility and platform
independence. We do not resolve this tension by hiding it; we describe
it precisely.

\paragraph{What Tinker is and is not.}
Tinker \cite{thinkingmachines2024tinker} is a managed LLM fine-tuning and
inference service provided by Thinking Machines, Inc. It exposes a Python
SDK (\texttt{tinker==0.16.1}) that accepts custom loss functions
(\texttt{forward\_backward\_custom}), a limited set of optimisers, and
standard LoRA hyperparameters. It does not expose (i) the exact server-side
GRPO loss implementation, (ii) reward normalisation or baseline subtraction
scheme, (iii) minibatch construction or gradient-accumulation strategy,
(iv) hardware configuration (GPU type, inter-node bandwidth), or (v)
system-level telemetry (energy, throughput, queueing).

\paragraph{What this means for our claims.}
Tinker results in this paper measure the \emph{platform's} implementation
of GRPO, not an abstract specification of the algorithm. A reader who asks
``why does Tinker GRPO score 99.9\% on GSM8K while open-source TRL GRPO
scores 73.4\% on the same task'' cannot fully answer that question from
our data: we are able to rule out a handful of candidate explanations
(seed variance, model-size confound) but not the implementation itself.
We therefore draw \emph{quantitative} conclusions only from the
open-source side (TRL, veRL on Modal H100, OpenRLHF) where every
hyperparameter is auditable, and use Tinker results primarily as an
\emph{upper bound} on what a carefully engineered production stack can
achieve for critic-free RL at our scales.

\paragraph{Reproducibility commitments we can make.}
\begin{enumerate}
  \item All Tinker experiment scripts, configuration files, JSON step
        logs, and per-run W\&B projects are archived in the repository.
        Researchers with Tinker access can attempt replication with our
        exact configurations.
  \item Every figure and every summary statistic derived solely from
        Tinker data is marked with ``\dag'' and a footnote stating that
        independent replication requires Tinker API access.
  \item Primary statistical claims (significance tests, ANOVA variance
        decompositions) are drawn from open-source Modal H100 and TRL
        runs. Tinker-only results are reported descriptively.
  \item We commit to re-running any Tinker-only experiment on an open
        backend if an equivalent open-source service becomes available,
        and to issuing a revised version of this paper if the Tinker
        99.9\% figure is ever shown to be due to an undisclosed
        implementation choice rather than the algorithm.
\end{enumerate}

\paragraph{Key rotation and credential hygiene.}
Our Tinker API key (prefix \texttt{tml-...}) is held in a repository
\texttt{.env.example} placeholder only. The real key is stored in the
authors' password manager and rotated whenever a failure mode suggests
possible exfiltration. No Tinker key is committed to git history in
either the primary repository (\texttt{pes-llm-research/tinker-rl-lab}) or
the mirror (\texttt{arvindcr4/tinker-rl-lab}).

% ---------------------------------------------------------------------------
\subsection{Known Methodological Limits}
\label{sec:methodlimits}
% ---------------------------------------------------------------------------

In addition to the infrastructure failures and platform-specific caveats
enumerated in Section~\ref{sec:limitations}, we flag the following
methodological limits.

\paragraph{Short training horizons.}
All Tinker GRPO runs were capped at 30--50 gradient steps. This is a
deliberate cost-control choice and means our Tinker numbers are
\emph{early-training snapshots}. We cannot rule out that longer training
would change the qualitative picture (e.g., the 82\%$\to$83.3\% held-out
gain might widen, or might collapse to reward hacking).

\paragraph{Single-seed Tinker experiments.}
Each Tinker configuration was run once, not 3--5 times as best practice
prescribes \cite{henderson2018deep}. Tinker results therefore lack
variance estimates and do not support significance testing; we report
them descriptively. Only the 5-seed TRL baseline and the 5-seed held-out
GSM8K evaluation carry proper confidence intervals.

\paragraph{Train-set reward as primary Tinker metric.}
Tinker runs primarily report reward on training prompts. Only the GSM8K
track was followed up with a separate 200-example held-out evaluation
per seed. Other Tinker tracks (tool-use, xLAM) remain train-set-only and
we cannot distinguish memorisation from generalisation for those
settings.

\paragraph{LoRA only; no full fine-tuning.}
All experiments use LoRA \cite{hu2022lora} with ranks 8--64. We have
\emph{not} tested whether full fine-tuning would re-order the
library/algorithm comparisons. The LoRA constraint is practical (cost) but
it means our claims about PPO vs.\ GRPO and TRL vs.\ Tinker are
LoRA-conditional.

\paragraph{Benchmark coverage is narrow.}
We evaluate four domains: GSM8K, MATH-500 (exploratory), HumanEval
(subset), and synthetic/xLAM tool-use. We do not evaluate on MT-Bench,
ArenaHard, safety benchmarks (HarmBench, ToxicChat), or truthfulness
benchmarks (TruthfulQA). Any claim about ``reasoning'' in the paper is
strictly a claim about the covered benchmarks.

\paragraph{No human preference data.}
We use verifiable rewards only (exact-match for math, unit-test for code,
schema-validity for tool-use). We do not study RLHF with human preference
data; findings should not be extrapolated to reward-model--based alignment
without further work.

\paragraph{Closed-source implementation opacity (reiterated).}
Comparisons between Tinker GRPO and TRL GRPO are confounded by the
closed-source nature of Tinker's implementation. See
Section~\ref{sec:tinker_closed} for detail.

\paragraph{Cross-platform hardware confounding.}
Tinker, Modal, Colab, and the PES A100 node use different accelerators,
different memory hierarchies, and different inference/training stacks
(Tinker proprietary, Modal vLLM, Colab PyTorch-native, TRL+Accelerate).
Observed differences in reward dynamics are not purely algorithmic.

\paragraph{Carbon estimates are approximate.}
As discussed in Section~\ref{sec:carbon}, Tinker GPU-hour and TDP are
inferred, not measured. Grid intensity is region-average, not marginal.
Numbers in Table~\ref{tab:carbon2} should be treated as
order-of-magnitude.

\paragraph{Exploratory, not confirmatory.}
This is an exploratory case study, not a pre-registered confirmatory
experiment. We did not pre-register primary hypotheses. All $p$-values are
un-corrected for the number of comparisons actually made across the
paper; the Bonferroni-surviving comparisons in
Section~\ref{sec:results} are so flagged, but most of our secondary
analyses are descriptive.

\paragraph{Geographic and demographic limits.}
All authors are based at a single institution in Bengaluru, India. Our
choice of benchmarks, prompt phrasings, and evaluation design reflects
that context. Independent replication by groups in other regions, on
non-English tasks, and with non-Qwen / non-Llama families is an open
question.

% =============================================================================
% End of ethics_statement.tex
 into main.tex
% Intended placement: after the Conclusion section, before References.
% Label: \label{sec:impact_standalone}
% Target length: >= 2 pages (NeurIPS Broader Impact guidance).
% =============================================================================

\section{Ethics, Limitations, and Broader Impact}
\label{sec:impact_standalone}

This statement is written to satisfy the NeurIPS Code of Ethics and Broader
Impact requirements. It complements the shorter Limitations and Broader
Impact subsections embedded in Section~\ref{sec:limitations} by consolidating
in one place a full dual-use analysis, itemised compute accounting, carbon
footprint estimation, data provenance disclosures, a candid acknowledgment
of the closed-source training backend on which a fraction of our headline
numbers depends, and a list of methodological limits we are aware of but
did not have resources to close within the submission window.

% ---------------------------------------------------------------------------
\subsection{Dual-Use Analysis: Misuse Risks of Reasoning RL}
\label{sec:dualuse}
% ---------------------------------------------------------------------------

\paragraph{Threat model.}
\ours{} releases (i) training scripts that apply GRPO to the GSM8K
\cite{cobbe2021training} mathematical reasoning benchmark and to the
Salesforce xLAM-Function-Calling-60k \cite{liu2024xlam} tool-use corpus,
(ii) a set of LoRA adapters published on the HuggingFace Hub
(\texttt{arvindcr4/tinker-rl-bench-*}), and (iii) diagnostic code for
Zero-Variance Fraction, reward-stability, and length-bias analysis. We
analyse four concrete misuse pathways these artefacts could, in principle,
enable, and we describe the existing guardrails and the residual risk we
judge acceptable for publication.

\paragraph{M1. Weaponisation of mathematical reasoning.}
The most capability-relevant artefact we release is GRPO training code that
improves grade-school arithmetic and multi-step reasoning. A malicious
operator could attempt to extend this pipeline to tasks of greater
dual-use concern, e.g., chemistry olympiad problems, cryptographic key
recovery, or financial fraud. We judge the \emph{marginal} uplift provided
by our release to be small for three reasons. First, GRPO at our training
scale does not create new capabilities; the base-model control for GSM8K
(Section~\ref{sec:results}, 82.0\%) shows that the bulk of held-out
accuracy is attributable to Qwen3-8B's pretraining. Our full GRPO run only
adds \(+1.3\) percentage points ($p{=}0.26$, not statistically significant).
Second, the public literature already contains GRPO implementations
\cite{shao2024deepseekmath, openrlhf2024, verl2024, trl2020} that are at
least as capable. Third, any reasoning uplift is bounded by the rewarder:
GSM8K rewards use exact-match against a human-written gold answer, which
does not generalise to the domains listed above without a new reward
signal, which itself requires expertise and labelled data.

\paragraph{M2. Reward hacking and the long-term alignment risk.}
GRPO, like all policy-gradient RL from a proxy reward, is vulnerable to
reward hacking: the model learns to exploit idiosyncrasies of the reward
function rather than solve the intended task \cite{skalse2022specificationgaming,
krakovna2020specification}. Our length-bias analysis
(Section~\ref{sec:length_bias}) and the verbosity-trap finding in 4/11
GRPO runs are concrete demonstrations of this failure mode in our own data.
Users who apply our scripts to under-specified reward functions in
production settings (e.g., a custom ``helpfulness'' classifier) may observe
models that appear to improve on held-out metrics while silently
optimising against unintended proxies. We include the Zero-Variance
Fraction and reward-stability diagnostics precisely so that downstream
users can detect such drift.

\paragraph{M3. Circumventing safety training via distillation or tool-use.}
The tool-use track trains LoRA adapters that teach a base model to emit
schema-valid JSON function calls \cite{liu2024xlam}. A motivated adversary
could in principle use the same pipeline to teach a base model to emit
jailbreak prompts as ``tool calls'' or to invoke an adversarial external
tool. We mitigate this by (a) releasing only adapters, not merged weights,
so that the safety-trained instruct checkpoint must be applied separately
and any instruct-layer guardrails remain active; (b) documenting in model
cards that the tool-use adapters are trained on a narrow schema (five
synthetic tools) and have not been safety-evaluated for open-ended function
calling; and (c) not releasing any function-calling harness that would
actually execute emitted calls.

\paragraph{M4. Compute-efficiency externalities.}
Our work argues that critic-free RL with careful group-size selection
(\(G{=}32\)) can be run on modest budgets. Improved training efficiency is
itself dual-use: it lowers the cost of running adversarial fine-tuning at
scale. We judge this externality to be modest relative to the already-low
cost of fine-tuning a 0.6B--8B model with commodity LoRA recipes
\cite{hu2022lora}, but we flag it explicitly.

\paragraph{Residual risk.}
After weighing M1--M4 we judge that the reproducibility, calibration, and
pedagogical benefits of releasing \ours{} outweigh the residual misuse
risk. We adopt the standard mitigation of releasing adapters (not merged
weights), documenting intended use and out-of-scope use in model cards,
and publishing alongside the diagnostics a researcher would need to detect
reward hacking.

% ---------------------------------------------------------------------------
\subsection{Compute Cost Accounting}
\label{sec:compute_accounting}
% ---------------------------------------------------------------------------

Table~\ref{tab:compute_spend} itemises the financial cost of every
platform used in the project. Costs are real spend for the submitting
authors; no in-kind credits are omitted. Dollar figures are in USD.

\begin{table}[ht]
\centering
\small
\caption{Itemised compute spend across platforms. Totals include both
successful and failed runs, because failed-run spend is real and should
not be hidden from reviewers \cite{henderson2018deep}.}
\label{tab:compute_spend}
\begin{tabular}{llrr}
\toprule
\textbf{Platform} & \textbf{Use} & \textbf{Runs} & \textbf{Cost (USD)} \\
\midrule
Tinker SDK v0.16.1 & GRPO on 4B/8B (original suite)        & 25 & \$40--55 \\
Tinker SDK v0.16.1 & 10x Structural Ceiling ablation       & 32 & \$65 \\
Tinker SDK v0.16.1 & World-Class Suite (frontier/MoE)      & 13 & \$25--30 \\
Modal H100         & PPO baselines (Qwen3-8B, Llama-3.1-8B)& 2  & \$12--18 \\
Modal H100         & KL-tracking run (failed)              & 1  & \$2--4   \\
Google Colab Pro (T4) & 0.5B--3B QLoRA SFT+GRPO            & --- & \$10/person \\
NVIDIA L4 (TRL baseline)& Qwen2.5-0.5B GSM8K, 5 seeds      & 5  & \$3--5   \\
HuggingFace Hub    & Model + adapter hosting               & --- & \$0 \\
Weights \& Biases  & Experiment tracking (academic tier)   & --- & \$0 \\
\midrule
\multicolumn{3}{l}{\textbf{Total authors' out-of-pocket spend}} & \textbf{\$130--140} \\
\bottomrule
\end{tabular}
\end{table}

\paragraph{Hidden costs.} Beyond direct platform spend, the project
consumed human labour (approximately six student-FTE-months, unpaid) and
infrastructure generously subsidised by PES University's LLM Research
Group (shared access to one A100 node used for TRL baselines and
pre-submission sanity checks). No author received commercial compensation
from Thinking Machines, Modal, or any other vendor mentioned.

\paragraph{Failed-run accounting.} Of the Tinker spend, we estimate
\textasciitilde\$30--35 was consumed by runs that ultimately failed
(JWT expiry, API stalls for GPT-OSS-20b and Kimi-K2 before re-run,
KL-tracking gradient bug). We disclose this because reviewers often ask
whether reported total spend reflects only clean, successful runs; it does
not. Excluding failed spend would understate the true resource cost of
reproducing our work by roughly 25\%.

% ---------------------------------------------------------------------------
\subsection{Carbon Footprint Estimation}
\label{sec:carbon}
% ---------------------------------------------------------------------------

Table~\ref{tab:carbon2} computes an energy and CO$_2$-equivalent estimate
for each platform. We follow the methodology of
\citet{patterson2024empirical} and \citet{strubell2019energy}: for each
GPU-hour we multiply device TDP by wall-clock GPU-hours and the data-centre
Power Usage Effectiveness (PUE), then multiply by the grid carbon intensity
in the region where the hardware is believed to be located.

\paragraph{Assumed constants.}
\begin{itemize}
  \item NVIDIA H100 SXM5 TDP: 700\,W.
  \item NVIDIA A100 80GB TDP: 400\,W.
  \item NVIDIA L4 TDP: 72\,W.
  \item NVIDIA T4 TDP: 70\,W.
  \item Data-centre PUE: 1.1 (conservative; typical hyperscale PUE is
        1.10--1.15 \cite{google2023pue, aws2023pue}).
  \item US average grid intensity (EPA eGRID 2023): 0.367\,kg\,CO$_2$-eq / kWh
        \cite{epa2023egrid}.
  \item Indian average grid intensity (CEA 2023): 0.716\,kg\,CO$_2$-eq / kWh
        \cite{cea2023co2}.
  \item We use the US average for Modal H100 (US-East region as configured)
        and Tinker (location undisclosed; assumed US), and the Indian
        average for Colab Pro T4 used by authors based in Bengaluru
        (Google Cloud asia-south1).
\end{itemize}

\paragraph{GPU-hour estimates.}
Tinker GPU-hour counts are not disclosed by the platform. We estimate them
from wall-clock run duration and assume a single H100-class accelerator per
job, which is consistent with Tinker's published architecture for the
model sizes we trained. Modal GPU-hours are measured directly from Modal's
billing dashboard.

\begin{table}[ht]
\centering
\small
\caption{Energy and carbon footprint estimates. Tinker hardware is not
publicly disclosed; we assume 1$\times$ H100 equivalent per run. Failed /
stalled runs are included. Totals rounded to one significant figure.}
\label{tab:carbon2}
\begin{tabular}{lrrrrr}
\toprule
\textbf{Platform} & \textbf{GPU-h} & \textbf{TDP\,(W)} & \textbf{PUE} &
\textbf{Energy\,(kWh)} & \textbf{CO$_2$\,(kg)} \\
\midrule
Tinker (assumed H100)   & \textasciitilde 950 & 700 & 1.1 & \textasciitilde 732 & \textasciitilde 269 \\
Modal H100 (US-East)    & \textasciitilde 18  & 700 & 1.1 & \textasciitilde 14  & \textasciitilde 5.1 \\
PES A100 (in-kind)      & \textasciitilde 60  & 400 & 1.1 & \textasciitilde 26  & \textasciitilde 19   \\
Colab Pro T4 (asia-south1) & \textasciitilde 40 & 70 & 1.2 & \textasciitilde 3.4 & \textasciitilde 2.4  \\
NVIDIA L4 (TRL baseline)   & \textasciitilde 10 & 72 & 1.1 & \textasciitilde 0.8 & \textasciitilde 0.3  \\
\midrule
\textbf{Project total (best estimate)} & & & & \textasciitilde 776 & \textasciitilde 296 \\
\bottomrule
\end{tabular}
\end{table}

\paragraph{Interpretation and uncertainty.}
Our best estimate of \textasciitilde 296\,kg\,CO$_2$-eq for the entire
project is comparable to a single round-trip economy flight
Bengaluru--Delhi (\textasciitilde 300\,kg). The dominant uncertainty is
the Tinker GPU-hour count: because Tinker does not expose hardware telemetry,
a factor-of-two error in GPU-hours or TDP is plausible. Under a pessimistic
assumption (2$\times$ GPU-hours, H100 running near 100\% TDP continuously),
the Tinker contribution could be as high as \textasciitilde 540\,kg CO$_2$-eq.
Under an optimistic assumption (Tinker backends actually use more
energy-efficient accelerators than H100, 50\% average utilisation) the
contribution could be as low as \textasciitilde 130\,kg. We report the
central estimate in the main table and caution readers that it should not
be interpreted as precise. \emph{All} numbers should be regarded as upper
bounds on the carbon cost of reproducing our work, because subsequent
users can skip the failed runs and the ablation sweeps.

\paragraph{Offset and mitigation.}
We did not purchase carbon offsets. Instead, we have (a) released all
trained checkpoints on HuggingFace Hub so that downstream users do not
need to re-run our experiments; (b) released step-level CSV logs so that
learning curves can be inspected without re-training; and (c) documented
in Section~\ref{sec:compute_accounting} which runs were wasteful, so that
replicators can skip them. We regard reproducibility-by-artefact as a
more durable mitigation than offsets.

% ---------------------------------------------------------------------------
\subsection{Data Provenance}
\label{sec:provenance}
% ---------------------------------------------------------------------------

\ours{} uses only publicly released research datasets. No private data,
personally identifiable information, or licensed proprietary content is
used. No human annotators were employed during this work. We document
each dataset below.

\paragraph{GSM8K \cite{cobbe2021training}.}
8{,}500 grade-school math word problems (7{,}473 train / 1{,}319 test)
authored by human writers contracted by OpenAI. Released under the
\textbf{MIT License} via \url{https://github.com/openai/grade-school-math}.
We use the standard train/test splits without modification. The canonical
citation is \citet{cobbe2021training}. Known limitations of GSM8K include
gender-neutral but culturally US-centric word problems (names, currencies,
sports) and a moderate rate of human-labelling errors estimated at
\textasciitilde 2\% by \cite{zhang2024gsm1k}. Our held-out evaluation
respects the test/train split.

\paragraph{Salesforce xLAM-Function-Calling-60k \cite{liu2024xlam}.}
60{,}000 function-calling examples generated by Salesforce Research as
training data for the xLAM agentic model family. Released under the
\textbf{CC BY 4.0} license via \url{https://huggingface.co/datasets/Salesforce/xlam-function-calling-60k}.
We use the public release as-is; specifically, we use the first 35 prompts
$\times$ 10 rollouts in the 10x Structural Ceiling tool-use track and the
full 60k split for the xlam-60k real-data run in Section~\ref{sec:results}.
Known limitations: xLAM schemas are synthetic and skew toward cleanly-typed
arguments; the dataset under-represents ambiguous, error-handling, and
multi-turn tool-calling. Attribution to Salesforce is required by the
license and is provided in Section~\ref{sec:acknowledgments} and in the
xLAM-derived model cards.

\paragraph{HumanEval \cite{chen2021codex}.}
164 Python programming problems with unit tests, released by OpenAI under
the \textbf{MIT License}. Used unmodified for pass@$k$ evaluation. Known
limitations: small scale, English-language docstrings, single-language
coverage; documented contamination concerns against frontier pretraining
corpora \cite{riddell2024contamination}.

\paragraph{NuminaMath \cite{li2024numinamath}.}
Used by collaborator Mohammad Rafi for the multi-stage
GSM8K+NuminaMath pipeline. Released under \textbf{Apache 2.0}. We use a
downstream checkpoint reported by Rafi; our repository contains only his
scripts and the Apache-licensed derivative weights.

\paragraph{Open-Platypus \cite{lee2023platypus}.}
3{,}000 SFT examples used by collaborator Madhu for Qwen3-8B code
generation warm-up. Open-Platypus is a curated subset released under
\textbf{CC BY-NC 4.0}; the non-commercial restriction is respected in our
release, which is academic-use only.

\paragraph{Synthetic tool-use corpus (authored).}
We additionally generated a small five-tool synthetic corpus (calculator,
web-search stub, calendar, file-read, email-send) used for Section 4.1's
saturation analysis. The corpus is authored by the paper's authors from
publicly documented APIs, contains no third-party content, and is
released under the \textbf{MIT License} in our repository under
\texttt{data/synthetic\_tools/}. Generation prompts are included for
auditability. No user data, scraped web content, or proprietary API
traces are present.

\paragraph{No web scraping, no human subjects.}
We did not scrape any website. We did not run any study with human
participants. No IRB review was therefore required. No data that could
reasonably be considered to contain PII, copyrighted content, or sensitive
personal attributes was used at any stage of training, evaluation, or
reward modelling.

% ---------------------------------------------------------------------------
\subsection{Closed-Source Tinker Acknowledgment}
\label{sec:tinker_closed}
% ---------------------------------------------------------------------------

A central tension in \ours{} is that a substantial fraction of our
headline numbers come from a closed-source commercial platform while our
paper simultaneously advocates for reproducibility and platform
independence. We do not resolve this tension by hiding it; we describe
it precisely.

\paragraph{What Tinker is and is not.}
Tinker \cite{thinkingmachines2024tinker} is a managed LLM fine-tuning and
inference service provided by Thinking Machines, Inc. It exposes a Python
SDK (\texttt{tinker==0.16.1}) that accepts custom loss functions
(\texttt{forward\_backward\_custom}), a limited set of optimisers, and
standard LoRA hyperparameters. It does not expose (i) the exact server-side
GRPO loss implementation, (ii) reward normalisation or baseline subtraction
scheme, (iii) minibatch construction or gradient-accumulation strategy,
(iv) hardware configuration (GPU type, inter-node bandwidth), or (v)
system-level telemetry (energy, throughput, queueing).

\paragraph{What this means for our claims.}
Tinker results in this paper measure the \emph{platform's} implementation
of GRPO, not an abstract specification of the algorithm. A reader who asks
``why does Tinker GRPO score 99.9\% on GSM8K while open-source TRL GRPO
scores 73.4\% on the same task'' cannot fully answer that question from
our data: we are able to rule out a handful of candidate explanations
(seed variance, model-size confound) but not the implementation itself.
We therefore draw \emph{quantitative} conclusions only from the
open-source side (TRL, veRL on Modal H100, OpenRLHF) where every
hyperparameter is auditable, and use Tinker results primarily as an
\emph{upper bound} on what a carefully engineered production stack can
achieve for critic-free RL at our scales.

\paragraph{Reproducibility commitments we can make.}
\begin{enumerate}
  \item All Tinker experiment scripts, configuration files, JSON step
        logs, and per-run W\&B projects are archived in the repository.
        Researchers with Tinker access can attempt replication with our
        exact configurations.
  \item Every figure and every summary statistic derived solely from
        Tinker data is marked with ``\dag'' and a footnote stating that
        independent replication requires Tinker API access.
  \item Primary statistical claims (significance tests, ANOVA variance
        decompositions) are drawn from open-source Modal H100 and TRL
        runs. Tinker-only results are reported descriptively.
  \item We commit to re-running any Tinker-only experiment on an open
        backend if an equivalent open-source service becomes available,
        and to issuing a revised version of this paper if the Tinker
        99.9\% figure is ever shown to be due to an undisclosed
        implementation choice rather than the algorithm.
\end{enumerate}

\paragraph{Key rotation and credential hygiene.}
Our Tinker API key (prefix \texttt{tml-...}) is held in a repository
\texttt{.env.example} placeholder only. The real key is stored in the
authors' password manager and rotated whenever a failure mode suggests
possible exfiltration. No Tinker key is committed to git history in
either the primary repository (\texttt{pes-llm-research/tinker-rl-lab}) or
the mirror (\texttt{arvindcr4/tinker-rl-lab}).

% ---------------------------------------------------------------------------
\subsection{Known Methodological Limits}
\label{sec:methodlimits}
% ---------------------------------------------------------------------------

In addition to the infrastructure failures and platform-specific caveats
enumerated in Section~\ref{sec:limitations}, we flag the following
methodological limits.

\paragraph{Short training horizons.}
All Tinker GRPO runs were capped at 30--50 gradient steps. This is a
deliberate cost-control choice and means our Tinker numbers are
\emph{early-training snapshots}. We cannot rule out that longer training
would change the qualitative picture (e.g., the 82\%$\to$83.3\% held-out
gain might widen, or might collapse to reward hacking).

\paragraph{Single-seed Tinker experiments.}
Each Tinker configuration was run once, not 3--5 times as best practice
prescribes \cite{henderson2018deep}. Tinker results therefore lack
variance estimates and do not support significance testing; we report
them descriptively. Only the 5-seed TRL baseline and the 5-seed held-out
GSM8K evaluation carry proper confidence intervals.

\paragraph{Train-set reward as primary Tinker metric.}
Tinker runs primarily report reward on training prompts. Only the GSM8K
track was followed up with a separate 200-example held-out evaluation
per seed. Other Tinker tracks (tool-use, xLAM) remain train-set-only and
we cannot distinguish memorisation from generalisation for those
settings.

\paragraph{LoRA only; no full fine-tuning.}
All experiments use LoRA \cite{hu2022lora} with ranks 8--64. We have
\emph{not} tested whether full fine-tuning would re-order the
library/algorithm comparisons. The LoRA constraint is practical (cost) but
it means our claims about PPO vs.\ GRPO and TRL vs.\ Tinker are
LoRA-conditional.

\paragraph{Benchmark coverage is narrow.}
We evaluate four domains: GSM8K, MATH-500 (exploratory), HumanEval
(subset), and synthetic/xLAM tool-use. We do not evaluate on MT-Bench,
ArenaHard, safety benchmarks (HarmBench, ToxicChat), or truthfulness
benchmarks (TruthfulQA). Any claim about ``reasoning'' in the paper is
strictly a claim about the covered benchmarks.

\paragraph{No human preference data.}
We use verifiable rewards only (exact-match for math, unit-test for code,
schema-validity for tool-use). We do not study RLHF with human preference
data; findings should not be extrapolated to reward-model--based alignment
without further work.

\paragraph{Closed-source implementation opacity (reiterated).}
Comparisons between Tinker GRPO and TRL GRPO are confounded by the
closed-source nature of Tinker's implementation. See
Section~\ref{sec:tinker_closed} for detail.

\paragraph{Cross-platform hardware confounding.}
Tinker, Modal, Colab, and the PES A100 node use different accelerators,
different memory hierarchies, and different inference/training stacks
(Tinker proprietary, Modal vLLM, Colab PyTorch-native, TRL+Accelerate).
Observed differences in reward dynamics are not purely algorithmic.

\paragraph{Carbon estimates are approximate.}
As discussed in Section~\ref{sec:carbon}, Tinker GPU-hour and TDP are
inferred, not measured. Grid intensity is region-average, not marginal.
Numbers in Table~\ref{tab:carbon2} should be treated as
order-of-magnitude.

\paragraph{Exploratory, not confirmatory.}
This is an exploratory case study, not a pre-registered confirmatory
experiment. We did not pre-register primary hypotheses. All $p$-values are
un-corrected for the number of comparisons actually made across the
paper; the Bonferroni-surviving comparisons in
Section~\ref{sec:results} are so flagged, but most of our secondary
analyses are descriptive.

\paragraph{Geographic and demographic limits.}
All authors are based at a single institution in Bengaluru, India. Our
choice of benchmarks, prompt phrasings, and evaluation design reflects
that context. Independent replication by groups in other regions, on
non-English tasks, and with non-Qwen / non-Llama families is an open
question.

% =============================================================================
% End of ethics_statement.tex
 into main.tex
% Intended placement: after the Conclusion section, before References.
% Label: \label{sec:impact_standalone}
% Target length: >= 2 pages (NeurIPS Broader Impact guidance).
% =============================================================================

\section{Ethics, Limitations, and Broader Impact}
\label{sec:impact_standalone}

This statement is written to satisfy the NeurIPS Code of Ethics and Broader
Impact requirements. It complements the shorter Limitations and Broader
Impact subsections embedded in Section~\ref{sec:limitations} by consolidating
in one place a full dual-use analysis, itemised compute accounting, carbon
footprint estimation, data provenance disclosures, a candid acknowledgment
of the closed-source training backend on which a fraction of our headline
numbers depends, and a list of methodological limits we are aware of but
did not have resources to close within the submission window.

% ---------------------------------------------------------------------------
\subsection{Dual-Use Analysis: Misuse Risks of Reasoning RL}
\label{sec:dualuse}
% ---------------------------------------------------------------------------

\paragraph{Threat model.}
\ours{} releases (i) training scripts that apply GRPO to the GSM8K
\cite{cobbe2021training} mathematical reasoning benchmark and to the
Salesforce xLAM-Function-Calling-60k \cite{liu2024xlam} tool-use corpus,
(ii) a set of LoRA adapters published on the HuggingFace Hub
(\texttt{arvindcr4/tinker-rl-bench-*}), and (iii) diagnostic code for
Zero-Variance Fraction, reward-stability, and length-bias analysis. We
analyse four concrete misuse pathways these artefacts could, in principle,
enable, and we describe the existing guardrails and the residual risk we
judge acceptable for publication.

\paragraph{M1. Weaponisation of mathematical reasoning.}
The most capability-relevant artefact we release is GRPO training code that
improves grade-school arithmetic and multi-step reasoning. A malicious
operator could attempt to extend this pipeline to tasks of greater
dual-use concern, e.g., chemistry olympiad problems, cryptographic key
recovery, or financial fraud. We judge the \emph{marginal} uplift provided
by our release to be small for three reasons. First, GRPO at our training
scale does not create new capabilities; the base-model control for GSM8K
(Section~\ref{sec:results}, 82.0\%) shows that the bulk of held-out
accuracy is attributable to Qwen3-8B's pretraining. Our full GRPO run only
adds \(+1.3\) percentage points ($p{=}0.26$, not statistically significant).
Second, the public literature already contains GRPO implementations
\cite{shao2024deepseekmath, openrlhf2024, verl2024, trl2020} that are at
least as capable. Third, any reasoning uplift is bounded by the rewarder:
GSM8K rewards use exact-match against a human-written gold answer, which
does not generalise to the domains listed above without a new reward
signal, which itself requires expertise and labelled data.

\paragraph{M2. Reward hacking and the long-term alignment risk.}
GRPO, like all policy-gradient RL from a proxy reward, is vulnerable to
reward hacking: the model learns to exploit idiosyncrasies of the reward
function rather than solve the intended task \cite{skalse2022specificationgaming,
krakovna2020specification}. Our length-bias analysis
(Section~\ref{sec:length_bias}) and the verbosity-trap finding in 4/11
GRPO runs are concrete demonstrations of this failure mode in our own data.
Users who apply our scripts to under-specified reward functions in
production settings (e.g., a custom ``helpfulness'' classifier) may observe
models that appear to improve on held-out metrics while silently
optimising against unintended proxies. We include the Zero-Variance
Fraction and reward-stability diagnostics precisely so that downstream
users can detect such drift.

\paragraph{M3. Circumventing safety training via distillation or tool-use.}
The tool-use track trains LoRA adapters that teach a base model to emit
schema-valid JSON function calls \cite{liu2024xlam}. A motivated adversary
could in principle use the same pipeline to teach a base model to emit
jailbreak prompts as ``tool calls'' or to invoke an adversarial external
tool. We mitigate this by (a) releasing only adapters, not merged weights,
so that the safety-trained instruct checkpoint must be applied separately
and any instruct-layer guardrails remain active; (b) documenting in model
cards that the tool-use adapters are trained on a narrow schema (five
synthetic tools) and have not been safety-evaluated for open-ended function
calling; and (c) not releasing any function-calling harness that would
actually execute emitted calls.

\paragraph{M4. Compute-efficiency externalities.}
Our work argues that critic-free RL with careful group-size selection
(\(G{=}32\)) can be run on modest budgets. Improved training efficiency is
itself dual-use: it lowers the cost of running adversarial fine-tuning at
scale. We judge this externality to be modest relative to the already-low
cost of fine-tuning a 0.6B--8B model with commodity LoRA recipes
\cite{hu2022lora}, but we flag it explicitly.

\paragraph{Residual risk.}
After weighing M1--M4 we judge that the reproducibility, calibration, and
pedagogical benefits of releasing \ours{} outweigh the residual misuse
risk. We adopt the standard mitigation of releasing adapters (not merged
weights), documenting intended use and out-of-scope use in model cards,
and publishing alongside the diagnostics a researcher would need to detect
reward hacking.

% ---------------------------------------------------------------------------
\subsection{Compute Cost Accounting}
\label{sec:compute_accounting}
% ---------------------------------------------------------------------------

Table~\ref{tab:compute_spend} itemises the financial cost of every
platform used in the project. Costs are real spend for the submitting
authors; no in-kind credits are omitted. Dollar figures are in USD.

\begin{table}[ht]
\centering
\small
\caption{Itemised compute spend across platforms. Totals include both
successful and failed runs, because failed-run spend is real and should
not be hidden from reviewers \cite{henderson2018deep}.}
\label{tab:compute_spend}
\begin{tabular}{llrr}
\toprule
\textbf{Platform} & \textbf{Use} & \textbf{Runs} & \textbf{Cost (USD)} \\
\midrule
Tinker SDK v0.16.1 & GRPO on 4B/8B (original suite)        & 25 & \$40--55 \\
Tinker SDK v0.16.1 & 10x Structural Ceiling ablation       & 32 & \$65 \\
Tinker SDK v0.16.1 & World-Class Suite (frontier/MoE)      & 13 & \$25--30 \\
Modal H100         & PPO baselines (Qwen3-8B, Llama-3.1-8B)& 2  & \$12--18 \\
Modal H100         & KL-tracking run (failed)              & 1  & \$2--4   \\
Google Colab Pro (T4) & 0.5B--3B QLoRA SFT+GRPO            & --- & \$10/person \\
NVIDIA L4 (TRL baseline)& Qwen2.5-0.5B GSM8K, 5 seeds      & 5  & \$3--5   \\
HuggingFace Hub    & Model + adapter hosting               & --- & \$0 \\
Weights \& Biases  & Experiment tracking (academic tier)   & --- & \$0 \\
\midrule
\multicolumn{3}{l}{\textbf{Total authors' out-of-pocket spend}} & \textbf{\$130--140} \\
\bottomrule
\end{tabular}
\end{table}

\paragraph{Hidden costs.} Beyond direct platform spend, the project
consumed human labour (approximately six student-FTE-months, unpaid) and
infrastructure generously subsidised by PES University's LLM Research
Group (shared access to one A100 node used for TRL baselines and
pre-submission sanity checks). No author received commercial compensation
from Thinking Machines, Modal, or any other vendor mentioned.

\paragraph{Failed-run accounting.} Of the Tinker spend, we estimate
\textasciitilde\$30--35 was consumed by runs that ultimately failed
(JWT expiry, API stalls for GPT-OSS-20b and Kimi-K2 before re-run,
KL-tracking gradient bug). We disclose this because reviewers often ask
whether reported total spend reflects only clean, successful runs; it does
not. Excluding failed spend would understate the true resource cost of
reproducing our work by roughly 25\%.

% ---------------------------------------------------------------------------
\subsection{Carbon Footprint Estimation}
\label{sec:carbon}
% ---------------------------------------------------------------------------

Table~\ref{tab:carbon2} computes an energy and CO$_2$-equivalent estimate
for each platform. We follow the methodology of
\citet{patterson2024empirical} and \citet{strubell2019energy}: for each
GPU-hour we multiply device TDP by wall-clock GPU-hours and the data-centre
Power Usage Effectiveness (PUE), then multiply by the grid carbon intensity
in the region where the hardware is believed to be located.

\paragraph{Assumed constants.}
\begin{itemize}
  \item NVIDIA H100 SXM5 TDP: 700\,W.
  \item NVIDIA A100 80GB TDP: 400\,W.
  \item NVIDIA L4 TDP: 72\,W.
  \item NVIDIA T4 TDP: 70\,W.
  \item Data-centre PUE: 1.1 (conservative; typical hyperscale PUE is
        1.10--1.15 \cite{google2023pue, aws2023pue}).
  \item US average grid intensity (EPA eGRID 2023): 0.367\,kg\,CO$_2$-eq / kWh
        \cite{epa2023egrid}.
  \item Indian average grid intensity (CEA 2023): 0.716\,kg\,CO$_2$-eq / kWh
        \cite{cea2023co2}.
  \item We use the US average for Modal H100 (US-East region as configured)
        and Tinker (location undisclosed; assumed US), and the Indian
        average for Colab Pro T4 used by authors based in Bengaluru
        (Google Cloud asia-south1).
\end{itemize}

\paragraph{GPU-hour estimates.}
Tinker GPU-hour counts are not disclosed by the platform. We estimate them
from wall-clock run duration and assume a single H100-class accelerator per
job, which is consistent with Tinker's published architecture for the
model sizes we trained. Modal GPU-hours are measured directly from Modal's
billing dashboard.

\begin{table}[ht]
\centering
\small
\caption{Energy and carbon footprint estimates. Tinker hardware is not
publicly disclosed; we assume 1$\times$ H100 equivalent per run. Failed /
stalled runs are included. Totals rounded to one significant figure.}
\label{tab:carbon2}
\begin{tabular}{lrrrrr}
\toprule
\textbf{Platform} & \textbf{GPU-h} & \textbf{TDP\,(W)} & \textbf{PUE} &
\textbf{Energy\,(kWh)} & \textbf{CO$_2$\,(kg)} \\
\midrule
Tinker (assumed H100)   & \textasciitilde 950 & 700 & 1.1 & \textasciitilde 732 & \textasciitilde 269 \\
Modal H100 (US-East)    & \textasciitilde 18  & 700 & 1.1 & \textasciitilde 14  & \textasciitilde 5.1 \\
PES A100 (in-kind)      & \textasciitilde 60  & 400 & 1.1 & \textasciitilde 26  & \textasciitilde 19   \\
Colab Pro T4 (asia-south1) & \textasciitilde 40 & 70 & 1.2 & \textasciitilde 3.4 & \textasciitilde 2.4  \\
NVIDIA L4 (TRL baseline)   & \textasciitilde 10 & 72 & 1.1 & \textasciitilde 0.8 & \textasciitilde 0.3  \\
\midrule
\textbf{Project total (best estimate)} & & & & \textasciitilde 776 & \textasciitilde 296 \\
\bottomrule
\end{tabular}
\end{table}

\paragraph{Interpretation and uncertainty.}
Our best estimate of \textasciitilde 296\,kg\,CO$_2$-eq for the entire
project is comparable to a single round-trip economy flight
Bengaluru--Delhi (\textasciitilde 300\,kg). The dominant uncertainty is
the Tinker GPU-hour count: because Tinker does not expose hardware telemetry,
a factor-of-two error in GPU-hours or TDP is plausible. Under a pessimistic
assumption (2$\times$ GPU-hours, H100 running near 100\% TDP continuously),
the Tinker contribution could be as high as \textasciitilde 540\,kg CO$_2$-eq.
Under an optimistic assumption (Tinker backends actually use more
energy-efficient accelerators than H100, 50\% average utilisation) the
contribution could be as low as \textasciitilde 130\,kg. We report the
central estimate in the main table and caution readers that it should not
be interpreted as precise. \emph{All} numbers should be regarded as upper
bounds on the carbon cost of reproducing our work, because subsequent
users can skip the failed runs and the ablation sweeps.

\paragraph{Offset and mitigation.}
We did not purchase carbon offsets. Instead, we have (a) released all
trained checkpoints on HuggingFace Hub so that downstream users do not
need to re-run our experiments; (b) released step-level CSV logs so that
learning curves can be inspected without re-training; and (c) documented
in Section~\ref{sec:compute_accounting} which runs were wasteful, so that
replicators can skip them. We regard reproducibility-by-artefact as a
more durable mitigation than offsets.

% ---------------------------------------------------------------------------
\subsection{Data Provenance}
\label{sec:provenance}
% ---------------------------------------------------------------------------

\ours{} uses only publicly released research datasets. No private data,
personally identifiable information, or licensed proprietary content is
used. No human annotators were employed during this work. We document
each dataset below.

\paragraph{GSM8K \cite{cobbe2021training}.}
8{,}500 grade-school math word problems (7{,}473 train / 1{,}319 test)
authored by human writers contracted by OpenAI. Released under the
\textbf{MIT License} via \url{https://github.com/openai/grade-school-math}.
We use the standard train/test splits without modification. The canonical
citation is \citet{cobbe2021training}. Known limitations of GSM8K include
gender-neutral but culturally US-centric word problems (names, currencies,
sports) and a moderate rate of human-labelling errors estimated at
\textasciitilde 2\% by \cite{zhang2024gsm1k}. Our held-out evaluation
respects the test/train split.

\paragraph{Salesforce xLAM-Function-Calling-60k \cite{liu2024xlam}.}
60{,}000 function-calling examples generated by Salesforce Research as
training data for the xLAM agentic model family. Released under the
\textbf{CC BY 4.0} license via \url{https://huggingface.co/datasets/Salesforce/xlam-function-calling-60k}.
We use the public release as-is; specifically, we use the first 35 prompts
$\times$ 10 rollouts in the 10x Structural Ceiling tool-use track and the
full 60k split for the xlam-60k real-data run in Section~\ref{sec:results}.
Known limitations: xLAM schemas are synthetic and skew toward cleanly-typed
arguments; the dataset under-represents ambiguous, error-handling, and
multi-turn tool-calling. Attribution to Salesforce is required by the
license and is provided in Section~\ref{sec:acknowledgments} and in the
xLAM-derived model cards.

\paragraph{HumanEval \cite{chen2021codex}.}
164 Python programming problems with unit tests, released by OpenAI under
the \textbf{MIT License}. Used unmodified for pass@$k$ evaluation. Known
limitations: small scale, English-language docstrings, single-language
coverage; documented contamination concerns against frontier pretraining
corpora \cite{riddell2024contamination}.

\paragraph{NuminaMath \cite{li2024numinamath}.}
Used by collaborator Mohammad Rafi for the multi-stage
GSM8K+NuminaMath pipeline. Released under \textbf{Apache 2.0}. We use a
downstream checkpoint reported by Rafi; our repository contains only his
scripts and the Apache-licensed derivative weights.

\paragraph{Open-Platypus \cite{lee2023platypus}.}
3{,}000 SFT examples used by collaborator Madhu for Qwen3-8B code
generation warm-up. Open-Platypus is a curated subset released under
\textbf{CC BY-NC 4.0}; the non-commercial restriction is respected in our
release, which is academic-use only.

\paragraph{Synthetic tool-use corpus (authored).}
We additionally generated a small five-tool synthetic corpus (calculator,
web-search stub, calendar, file-read, email-send) used for Section 4.1's
saturation analysis. The corpus is authored by the paper's authors from
publicly documented APIs, contains no third-party content, and is
released under the \textbf{MIT License} in our repository under
\texttt{data/synthetic\_tools/}. Generation prompts are included for
auditability. No user data, scraped web content, or proprietary API
traces are present.

\paragraph{No web scraping, no human subjects.}
We did not scrape any website. We did not run any study with human
participants. No IRB review was therefore required. No data that could
reasonably be considered to contain PII, copyrighted content, or sensitive
personal attributes was used at any stage of training, evaluation, or
reward modelling.

% ---------------------------------------------------------------------------
\subsection{Closed-Source Tinker Acknowledgment}
\label{sec:tinker_closed}
% ---------------------------------------------------------------------------

A central tension in \ours{} is that a substantial fraction of our
headline numbers come from a closed-source commercial platform while our
paper simultaneously advocates for reproducibility and platform
independence. We do not resolve this tension by hiding it; we describe
it precisely.

\paragraph{What Tinker is and is not.}
Tinker \cite{thinkingmachines2024tinker} is a managed LLM fine-tuning and
inference service provided by Thinking Machines, Inc. It exposes a Python
SDK (\texttt{tinker==0.16.1}) that accepts custom loss functions
(\texttt{forward\_backward\_custom}), a limited set of optimisers, and
standard LoRA hyperparameters. It does not expose (i) the exact server-side
GRPO loss implementation, (ii) reward normalisation or baseline subtraction
scheme, (iii) minibatch construction or gradient-accumulation strategy,
(iv) hardware configuration (GPU type, inter-node bandwidth), or (v)
system-level telemetry (energy, throughput, queueing).

\paragraph{What this means for our claims.}
Tinker results in this paper measure the \emph{platform's} implementation
of GRPO, not an abstract specification of the algorithm. A reader who asks
``why does Tinker GRPO score 99.9\% on GSM8K while open-source TRL GRPO
scores 73.4\% on the same task'' cannot fully answer that question from
our data: we are able to rule out a handful of candidate explanations
(seed variance, model-size confound) but not the implementation itself.
We therefore draw \emph{quantitative} conclusions only from the
open-source side (TRL, veRL on Modal H100, OpenRLHF) where every
hyperparameter is auditable, and use Tinker results primarily as an
\emph{upper bound} on what a carefully engineered production stack can
achieve for critic-free RL at our scales.

\paragraph{Reproducibility commitments we can make.}
\begin{enumerate}
  \item All Tinker experiment scripts, configuration files, JSON step
        logs, and per-run W\&B projects are archived in the repository.
        Researchers with Tinker access can attempt replication with our
        exact configurations.
  \item Every figure and every summary statistic derived solely from
        Tinker data is marked with ``\dag'' and a footnote stating that
        independent replication requires Tinker API access.
  \item Primary statistical claims (significance tests, ANOVA variance
        decompositions) are drawn from open-source Modal H100 and TRL
        runs. Tinker-only results are reported descriptively.
  \item We commit to re-running any Tinker-only experiment on an open
        backend if an equivalent open-source service becomes available,
        and to issuing a revised version of this paper if the Tinker
        99.9\% figure is ever shown to be due to an undisclosed
        implementation choice rather than the algorithm.
\end{enumerate}

\paragraph{Key rotation and credential hygiene.}
Our Tinker API key (prefix \texttt{tml-...}) is held in a repository
\texttt{.env.example} placeholder only. The real key is stored in the
authors' password manager and rotated whenever a failure mode suggests
possible exfiltration. No Tinker key is committed to git history in
either the primary repository (\texttt{pes-llm-research/tinker-rl-lab}) or
the mirror (\texttt{arvindcr4/tinker-rl-lab}).

% ---------------------------------------------------------------------------
\subsection{Known Methodological Limits}
\label{sec:methodlimits}
% ---------------------------------------------------------------------------

In addition to the infrastructure failures and platform-specific caveats
enumerated in Section~\ref{sec:limitations}, we flag the following
methodological limits.

\paragraph{Short training horizons.}
All Tinker GRPO runs were capped at 30--50 gradient steps. This is a
deliberate cost-control choice and means our Tinker numbers are
\emph{early-training snapshots}. We cannot rule out that longer training
would change the qualitative picture (e.g., the 82\%$\to$83.3\% held-out
gain might widen, or might collapse to reward hacking).

\paragraph{Single-seed Tinker experiments.}
Each Tinker configuration was run once, not 3--5 times as best practice
prescribes \cite{henderson2018deep}. Tinker results therefore lack
variance estimates and do not support significance testing; we report
them descriptively. Only the 5-seed TRL baseline and the 5-seed held-out
GSM8K evaluation carry proper confidence intervals.

\paragraph{Train-set reward as primary Tinker metric.}
Tinker runs primarily report reward on training prompts. Only the GSM8K
track was followed up with a separate 200-example held-out evaluation
per seed. Other Tinker tracks (tool-use, xLAM) remain train-set-only and
we cannot distinguish memorisation from generalisation for those
settings.

\paragraph{LoRA only; no full fine-tuning.}
All experiments use LoRA \cite{hu2022lora} with ranks 8--64. We have
\emph{not} tested whether full fine-tuning would re-order the
library/algorithm comparisons. The LoRA constraint is practical (cost) but
it means our claims about PPO vs.\ GRPO and TRL vs.\ Tinker are
LoRA-conditional.

\paragraph{Benchmark coverage is narrow.}
We evaluate four domains: GSM8K, MATH-500 (exploratory), HumanEval
(subset), and synthetic/xLAM tool-use. We do not evaluate on MT-Bench,
ArenaHard, safety benchmarks (HarmBench, ToxicChat), or truthfulness
benchmarks (TruthfulQA). Any claim about ``reasoning'' in the paper is
strictly a claim about the covered benchmarks.

\paragraph{No human preference data.}
We use verifiable rewards only (exact-match for math, unit-test for code,
schema-validity for tool-use). We do not study RLHF with human preference
data; findings should not be extrapolated to reward-model--based alignment
without further work.

\paragraph{Closed-source implementation opacity (reiterated).}
Comparisons between Tinker GRPO and TRL GRPO are confounded by the
closed-source nature of Tinker's implementation. See
Section~\ref{sec:tinker_closed} for detail.

\paragraph{Cross-platform hardware confounding.}
Tinker, Modal, Colab, and the PES A100 node use different accelerators,
different memory hierarchies, and different inference/training stacks
(Tinker proprietary, Modal vLLM, Colab PyTorch-native, TRL+Accelerate).
Observed differences in reward dynamics are not purely algorithmic.

\paragraph{Carbon estimates are approximate.}
As discussed in Section~\ref{sec:carbon}, Tinker GPU-hour and TDP are
inferred, not measured. Grid intensity is region-average, not marginal.
Numbers in Table~\ref{tab:carbon2} should be treated as
order-of-magnitude.

\paragraph{Exploratory, not confirmatory.}
This is an exploratory case study, not a pre-registered confirmatory
experiment. We did not pre-register primary hypotheses. All $p$-values are
un-corrected for the number of comparisons actually made across the
paper; the Bonferroni-surviving comparisons in
Section~\ref{sec:results} are so flagged, but most of our secondary
analyses are descriptive.

\paragraph{Geographic and demographic limits.}
All authors are based at a single institution in Bengaluru, India. Our
choice of benchmarks, prompt phrasings, and evaluation design reflects
that context. Independent replication by groups in other regions, on
non-English tasks, and with non-Qwen / non-Llama families is an open
question.

% =============================================================================
% End of ethics_statement.tex
 into main.tex
% Intended placement: after the Conclusion section, before References.
% Label: \label{sec:impact_standalone}
% Target length: >= 2 pages (NeurIPS Broader Impact guidance).
% =============================================================================

\section{Ethics, Limitations, and Broader Impact}
\label{sec:impact_standalone}

This statement is written to satisfy the NeurIPS Code of Ethics and Broader
Impact requirements. It complements the shorter Limitations and Broader
Impact subsections embedded in Section~\ref{sec:limitations} by consolidating
in one place a full dual-use analysis, itemised compute accounting, carbon
footprint estimation, data provenance disclosures, a candid acknowledgment
of the closed-source training backend on which a fraction of our headline
numbers depends, and a list of methodological limits we are aware of but
did not have resources to close within the submission window.

% ---------------------------------------------------------------------------
\subsection{Dual-Use Analysis: Misuse Risks of Reasoning RL}
\label{sec:dualuse}
% ---------------------------------------------------------------------------

\paragraph{Threat model.}
\ours{} releases (i) training scripts that apply GRPO to the GSM8K
\cite{cobbe2021training} mathematical reasoning benchmark and to the
Salesforce xLAM-Function-Calling-60k \cite{liu2024xlam} tool-use corpus,
(ii) a set of LoRA adapters published on the HuggingFace Hub
(\texttt{arvindcr4/tinker-rl-bench-*}), and (iii) diagnostic code for
Zero-Variance Fraction, reward-stability, and length-bias analysis. We
analyse four concrete misuse pathways these artefacts could, in principle,
enable, and we describe the existing guardrails and the residual risk we
judge acceptable for publication.

\paragraph{M1. Weaponisation of mathematical reasoning.}
The most capability-relevant artefact we release is GRPO training code that
improves grade-school arithmetic and multi-step reasoning. A malicious
operator could attempt to extend this pipeline to tasks of greater
dual-use concern, e.g., chemistry olympiad problems, cryptographic key
recovery, or financial fraud. We judge the \emph{marginal} uplift provided
by our release to be small for three reasons. First, GRPO at our training
scale does not create new capabilities; the base-model control for GSM8K
(Section~\ref{sec:results}, 82.0\%) shows that the bulk of held-out
accuracy is attributable to Qwen3-8B's pretraining. Our full GRPO run only
adds \(+1.3\) percentage points ($p{=}0.26$, not statistically significant).
Second, the public literature already contains GRPO implementations
\cite{shao2024deepseekmath, openrlhf2024, verl2024, trl2020} that are at
least as capable. Third, any reasoning uplift is bounded by the rewarder:
GSM8K rewards use exact-match against a human-written gold answer, which
does not generalise to the domains listed above without a new reward
signal, which itself requires expertise and labelled data.

\paragraph{M2. Reward hacking and the long-term alignment risk.}
GRPO, like all policy-gradient RL from a proxy reward, is vulnerable to
reward hacking: the model learns to exploit idiosyncrasies of the reward
function rather than solve the intended task \cite{skalse2022specificationgaming,
krakovna2020specification}. Our length-bias analysis
(Section~\ref{sec:length_bias}) and the verbosity-trap finding in 4/11
GRPO runs are concrete demonstrations of this failure mode in our own data.
Users who apply our scripts to under-specified reward functions in
production settings (e.g., a custom ``helpfulness'' classifier) may observe
models that appear to improve on held-out metrics while silently
optimising against unintended proxies. We include the Zero-Variance
Fraction and reward-stability diagnostics precisely so that downstream
users can detect such drift.

\paragraph{M3. Circumventing safety training via distillation or tool-use.}
The tool-use track trains LoRA adapters that teach a base model to emit
schema-valid JSON function calls \cite{liu2024xlam}. A motivated adversary
could in principle use the same pipeline to teach a base model to emit
jailbreak prompts as ``tool calls'' or to invoke an adversarial external
tool. We mitigate this by (a) releasing only adapters, not merged weights,
so that the safety-trained instruct checkpoint must be applied separately
and any instruct-layer guardrails remain active; (b) documenting in model
cards that the tool-use adapters are trained on a narrow schema (five
synthetic tools) and have not been safety-evaluated for open-ended function
calling; and (c) not releasing any function-calling harness that would
actually execute emitted calls.

\paragraph{M4. Compute-efficiency externalities.}
Our work argues that critic-free RL with careful group-size selection
(\(G{=}32\)) can be run on modest budgets. Improved training efficiency is
itself dual-use: it lowers the cost of running adversarial fine-tuning at
scale. We judge this externality to be modest relative to the already-low
cost of fine-tuning a 0.6B--8B model with commodity LoRA recipes
\cite{hu2022lora}, but we flag it explicitly.

\paragraph{Residual risk.}
After weighing M1--M4 we judge that the reproducibility, calibration, and
pedagogical benefits of releasing \ours{} outweigh the residual misuse
risk. We adopt the standard mitigation of releasing adapters (not merged
weights), documenting intended use and out-of-scope use in model cards,
and publishing alongside the diagnostics a researcher would need to detect
reward hacking.

% ---------------------------------------------------------------------------
\subsection{Compute Cost Accounting}
\label{sec:compute_accounting}
% ---------------------------------------------------------------------------

Table~\ref{tab:compute_spend} itemises the financial cost of every
platform used in the project. Costs are real spend for the submitting
authors; no in-kind credits are omitted. Dollar figures are in USD.

\begin{table}[ht]
\centering
\small
\caption{Itemised compute spend across platforms. Totals include both
successful and failed runs, because failed-run spend is real and should
not be hidden from reviewers \cite{henderson2018deep}.}
\label{tab:compute_spend}
\begin{tabular}{llrr}
\toprule
\textbf{Platform} & \textbf{Use} & \textbf{Runs} & \textbf{Cost (USD)} \\
\midrule
Tinker SDK v0.16.1 & GRPO on 4B/8B (original suite)        & 25 & \$40--55 \\
Tinker SDK v0.16.1 & 10x Structural Ceiling ablation       & 32 & \$65 \\
Tinker SDK v0.16.1 & World-Class Suite (frontier/MoE)      & 13 & \$25--30 \\
Modal H100         & PPO baselines (Qwen3-8B, Llama-3.1-8B)& 2  & \$12--18 \\
Modal H100         & KL-tracking run (failed)              & 1  & \$2--4   \\
Google Colab Pro (T4) & 0.5B--3B QLoRA SFT+GRPO            & --- & \$10/person \\
NVIDIA L4 (TRL baseline)& Qwen2.5-0.5B GSM8K, 5 seeds      & 5  & \$3--5   \\
HuggingFace Hub    & Model + adapter hosting               & --- & \$0 \\
Weights \& Biases  & Experiment tracking (academic tier)   & --- & \$0 \\
\midrule
\multicolumn{3}{l}{\textbf{Total authors' out-of-pocket spend}} & \textbf{\$130--140} \\
\bottomrule
\end{tabular}
\end{table}

\paragraph{Hidden costs.} Beyond direct platform spend, the project
consumed human labour (approximately six student-FTE-months, unpaid) and
infrastructure generously subsidised by PES University's LLM Research
Group (shared access to one A100 node used for TRL baselines and
pre-submission sanity checks). No author received commercial compensation
from Thinking Machines, Modal, or any other vendor mentioned.

\paragraph{Failed-run accounting.} Of the Tinker spend, we estimate
\textasciitilde\$30--35 was consumed by runs that ultimately failed
(JWT expiry, API stalls for GPT-OSS-20b and Kimi-K2 before re-run,
KL-tracking gradient bug). We disclose this because reviewers often ask
whether reported total spend reflects only clean, successful runs; it does
not. Excluding failed spend would understate the true resource cost of
reproducing our work by roughly 25\%.

% ---------------------------------------------------------------------------
\subsection{Carbon Footprint Estimation}
\label{sec:carbon}
% ---------------------------------------------------------------------------

Table~\ref{tab:carbon2} computes an energy and CO$_2$-equivalent estimate
for each platform. We follow the methodology of
\citet{patterson2024empirical} and \citet{strubell2019energy}: for each
GPU-hour we multiply device TDP by wall-clock GPU-hours and the data-centre
Power Usage Effectiveness (PUE), then multiply by the grid carbon intensity
in the region where the hardware is believed to be located.

\paragraph{Assumed constants.}
\begin{itemize}
  \item NVIDIA H100 SXM5 TDP: 700\,W.
  \item NVIDIA A100 80GB TDP: 400\,W.
  \item NVIDIA L4 TDP: 72\,W.
  \item NVIDIA T4 TDP: 70\,W.
  \item Data-centre PUE: 1.1 (conservative; typical hyperscale PUE is
        1.10--1.15 \cite{google2023pue, aws2023pue}).
  \item US average grid intensity (EPA eGRID 2023): 0.367\,kg\,CO$_2$-eq / kWh
        \cite{epa2023egrid}.
  \item Indian average grid intensity (CEA 2023): 0.716\,kg\,CO$_2$-eq / kWh
        \cite{cea2023co2}.
  \item We use the US average for Modal H100 (US-East region as configured)
        and Tinker (location undisclosed; assumed US), and the Indian
        average for Colab Pro T4 used by authors based in Bengaluru
        (Google Cloud asia-south1).
\end{itemize}

\paragraph{GPU-hour estimates.}
Tinker GPU-hour counts are not disclosed by the platform. We estimate them
from wall-clock run duration and assume a single H100-class accelerator per
job, which is consistent with Tinker's published architecture for the
model sizes we trained. Modal GPU-hours are measured directly from Modal's
billing dashboard.

\begin{table}[ht]
\centering
\small
\caption{Energy and carbon footprint estimates. Tinker hardware is not
publicly disclosed; we assume 1$\times$ H100 equivalent per run. Failed /
stalled runs are included. Totals rounded to one significant figure.}
\label{tab:carbon2}
\begin{tabular}{lrrrrr}
\toprule
\textbf{Platform} & \textbf{GPU-h} & \textbf{TDP\,(W)} & \textbf{PUE} &
\textbf{Energy\,(kWh)} & \textbf{CO$_2$\,(kg)} \\
\midrule
Tinker (assumed H100)   & \textasciitilde 950 & 700 & 1.1 & \textasciitilde 732 & \textasciitilde 269 \\
Modal H100 (US-East)    & \textasciitilde 18  & 700 & 1.1 & \textasciitilde 14  & \textasciitilde 5.1 \\
PES A100 (in-kind)      & \textasciitilde 60  & 400 & 1.1 & \textasciitilde 26  & \textasciitilde 19   \\
Colab Pro T4 (asia-south1) & \textasciitilde 40 & 70 & 1.2 & \textasciitilde 3.4 & \textasciitilde 2.4  \\
NVIDIA L4 (TRL baseline)   & \textasciitilde 10 & 72 & 1.1 & \textasciitilde 0.8 & \textasciitilde 0.3  \\
\midrule
\textbf{Project total (best estimate)} & & & & \textasciitilde 776 & \textasciitilde 296 \\
\bottomrule
\end{tabular}
\end{table}

\paragraph{Interpretation and uncertainty.}
Our best estimate of \textasciitilde 296\,kg\,CO$_2$-eq for the entire
project is comparable to a single round-trip economy flight
Bengaluru--Delhi (\textasciitilde 300\,kg). The dominant uncertainty is
the Tinker GPU-hour count: because Tinker does not expose hardware telemetry,
a factor-of-two error in GPU-hours or TDP is plausible. Under a pessimistic
assumption (2$\times$ GPU-hours, H100 running near 100\% TDP continuously),
the Tinker contribution could be as high as \textasciitilde 540\,kg CO$_2$-eq.
Under an optimistic assumption (Tinker backends actually use more
energy-efficient accelerators than H100, 50\% average utilisation) the
contribution could be as low as \textasciitilde 130\,kg. We report the
central estimate in the main table and caution readers that it should not
be interpreted as precise. \emph{All} numbers should be regarded as upper
bounds on the carbon cost of reproducing our work, because subsequent
users can skip the failed runs and the ablation sweeps.

\paragraph{Offset and mitigation.}
We did not purchase carbon offsets. Instead, we have (a) released all
trained checkpoints on HuggingFace Hub so that downstream users do not
need to re-run our experiments; (b) released step-level CSV logs so that
learning curves can be inspected without re-training; and (c) documented
in Section~\ref{sec:compute_accounting} which runs were wasteful, so that
replicators can skip them. We regard reproducibility-by-artefact as a
more durable mitigation than offsets.

% ---------------------------------------------------------------------------
\subsection{Data Provenance}
\label{sec:provenance}
% ---------------------------------------------------------------------------

\ours{} uses only publicly released research datasets. No private data,
personally identifiable information, or licensed proprietary content is
used. No human annotators were employed during this work. We document
each dataset below.

\paragraph{GSM8K \cite{cobbe2021training}.}
8{,}500 grade-school math word problems (7{,}473 train / 1{,}319 test)
authored by human writers contracted by OpenAI. Released under the
\textbf{MIT License} via \url{https://github.com/openai/grade-school-math}.
We use the standard train/test splits without modification. The canonical
citation is \citet{cobbe2021training}. Known limitations of GSM8K include
gender-neutral but culturally US-centric word problems (names, currencies,
sports) and a moderate rate of human-labelling errors estimated at
\textasciitilde 2\% by \cite{zhang2024gsm1k}. Our held-out evaluation
respects the test/train split.

\paragraph{Salesforce xLAM-Function-Calling-60k \cite{liu2024xlam}.}
60{,}000 function-calling examples generated by Salesforce Research as
training data for the xLAM agentic model family. Released under the
\textbf{CC BY 4.0} license via \url{https://huggingface.co/datasets/Salesforce/xlam-function-calling-60k}.
We use the public release as-is; specifically, we use the first 35 prompts
$\times$ 10 rollouts in the 10x Structural Ceiling tool-use track and the
full 60k split for the xlam-60k real-data run in Section~\ref{sec:results}.
Known limitations: xLAM schemas are synthetic and skew toward cleanly-typed
arguments; the dataset under-represents ambiguous, error-handling, and
multi-turn tool-calling. Attribution to Salesforce is required by the
license and is provided in Section~\ref{sec:acknowledgments} and in the
xLAM-derived model cards.

\paragraph{HumanEval \cite{chen2021codex}.}
164 Python programming problems with unit tests, released by OpenAI under
the \textbf{MIT License}. Used unmodified for pass@$k$ evaluation. Known
limitations: small scale, English-language docstrings, single-language
coverage; documented contamination concerns against frontier pretraining
corpora \cite{riddell2024contamination}.

\paragraph{NuminaMath \cite{li2024numinamath}.}
Used by collaborator Mohammad Rafi for the multi-stage
GSM8K+NuminaMath pipeline. Released under \textbf{Apache 2.0}. We use a
downstream checkpoint reported by Rafi; our repository contains only his
scripts and the Apache-licensed derivative weights.

\paragraph{Open-Platypus \cite{lee2023platypus}.}
3{,}000 SFT examples used by collaborator Madhu for Qwen3-8B code
generation warm-up. Open-Platypus is a curated subset released under
\textbf{CC BY-NC 4.0}; the non-commercial restriction is respected in our
release, which is academic-use only.

\paragraph{Synthetic tool-use corpus (authored).}
We additionally generated a small five-tool synthetic corpus (calculator,
web-search stub, calendar, file-read, email-send) used for Section 4.1's
saturation analysis. The corpus is authored by the paper's authors from
publicly documented APIs, contains no third-party content, and is
released under the \textbf{MIT License} in our repository under
\texttt{data/synthetic\_tools/}. Generation prompts are included for
auditability. No user data, scraped web content, or proprietary API
traces are present.

\paragraph{No web scraping, no human subjects.}
We did not scrape any website. We did not run any study with human
participants. No IRB review was therefore required. No data that could
reasonably be considered to contain PII, copyrighted content, or sensitive
personal attributes was used at any stage of training, evaluation, or
reward modelling.

% ---------------------------------------------------------------------------
\subsection{Closed-Source Tinker Acknowledgment}
\label{sec:tinker_closed}
% ---------------------------------------------------------------------------

A central tension in \ours{} is that a substantial fraction of our
headline numbers come from a closed-source commercial platform while our
paper simultaneously advocates for reproducibility and platform
independence. We do not resolve this tension by hiding it; we describe
it precisely.

\paragraph{What Tinker is and is not.}
Tinker \cite{thinkingmachines2024tinker} is a managed LLM fine-tuning and
inference service provided by Thinking Machines, Inc. It exposes a Python
SDK (\texttt{tinker==0.16.1}) that accepts custom loss functions
(\texttt{forward\_backward\_custom}), a limited set of optimisers, and
standard LoRA hyperparameters. It does not expose (i) the exact server-side
GRPO loss implementation, (ii) reward normalisation or baseline subtraction
scheme, (iii) minibatch construction or gradient-accumulation strategy,
(iv) hardware configuration (GPU type, inter-node bandwidth), or (v)
system-level telemetry (energy, throughput, queueing).

\paragraph{What this means for our claims.}
Tinker results in this paper measure the \emph{platform's} implementation
of GRPO, not an abstract specification of the algorithm. A reader who asks
``why does Tinker GRPO score 99.9\% on GSM8K while open-source TRL GRPO
scores 73.4\% on the same task'' cannot fully answer that question from
our data: we are able to rule out a handful of candidate explanations
(seed variance, model-size confound) but not the implementation itself.
We therefore draw \emph{quantitative} conclusions only from the
open-source side (TRL, veRL on Modal H100, OpenRLHF) where every
hyperparameter is auditable, and use Tinker results primarily as an
\emph{upper bound} on what a carefully engineered production stack can
achieve for critic-free RL at our scales.

\paragraph{Reproducibility commitments we can make.}
\begin{enumerate}
  \item All Tinker experiment scripts, configuration files, JSON step
        logs, and per-run W\&B projects are archived in the repository.
        Researchers with Tinker access can attempt replication with our
        exact configurations.
  \item Every figure and every summary statistic derived solely from
        Tinker data is marked with ``\dag'' and a footnote stating that
        independent replication requires Tinker API access.
  \item Primary statistical claims (significance tests, ANOVA variance
        decompositions) are drawn from open-source Modal H100 and TRL
        runs. Tinker-only results are reported descriptively.
  \item We commit to re-running any Tinker-only experiment on an open
        backend if an equivalent open-source service becomes available,
        and to issuing a revised version of this paper if the Tinker
        99.9\% figure is ever shown to be due to an undisclosed
        implementation choice rather than the algorithm.
\end{enumerate}

\paragraph{Key rotation and credential hygiene.}
Our Tinker API key (prefix \texttt{tml-...}) is held in a repository
\texttt{.env.example} placeholder only. The real key is stored in the
authors' password manager and rotated whenever a failure mode suggests
possible exfiltration. No Tinker key is committed to git history in
either the primary repository (\texttt{pes-llm-research/tinker-rl-lab}) or
the mirror (\texttt{arvindcr4/tinker-rl-lab}).

% ---------------------------------------------------------------------------
\subsection{Known Methodological Limits}
\label{sec:methodlimits}
% ---------------------------------------------------------------------------

In addition to the infrastructure failures and platform-specific caveats
enumerated in Section~\ref{sec:limitations}, we flag the following
methodological limits.

\paragraph{Short training horizons.}
All Tinker GRPO runs were capped at 30--50 gradient steps. This is a
deliberate cost-control choice and means our Tinker numbers are
\emph{early-training snapshots}. We cannot rule out that longer training
would change the qualitative picture (e.g., the 82\%$\to$83.3\% held-out
gain might widen, or might collapse to reward hacking).

\paragraph{Single-seed Tinker experiments.}
Each Tinker configuration was run once, not 3--5 times as best practice
prescribes \cite{henderson2018deep}. Tinker results therefore lack
variance estimates and do not support significance testing; we report
them descriptively. Only the 5-seed TRL baseline and the 5-seed held-out
GSM8K evaluation carry proper confidence intervals.

\paragraph{Train-set reward as primary Tinker metric.}
Tinker runs primarily report reward on training prompts. Only the GSM8K
track was followed up with a separate 200-example held-out evaluation
per seed. Other Tinker tracks (tool-use, xLAM) remain train-set-only and
we cannot distinguish memorisation from generalisation for those
settings.

\paragraph{LoRA only; no full fine-tuning.}
All experiments use LoRA \cite{hu2022lora} with ranks 8--64. We have
\emph{not} tested whether full fine-tuning would re-order the
library/algorithm comparisons. The LoRA constraint is practical (cost) but
it means our claims about PPO vs.\ GRPO and TRL vs.\ Tinker are
LoRA-conditional.

\paragraph{Benchmark coverage is narrow.}
We evaluate four domains: GSM8K, MATH-500 (exploratory), HumanEval
(subset), and synthetic/xLAM tool-use. We do not evaluate on MT-Bench,
ArenaHard, safety benchmarks (HarmBench, ToxicChat), or truthfulness
benchmarks (TruthfulQA). Any claim about ``reasoning'' in the paper is
strictly a claim about the covered benchmarks.

\paragraph{No human preference data.}
We use verifiable rewards only (exact-match for math, unit-test for code,
schema-validity for tool-use). We do not study RLHF with human preference
data; findings should not be extrapolated to reward-model--based alignment
without further work.

\paragraph{Closed-source implementation opacity (reiterated).}
Comparisons between Tinker GRPO and TRL GRPO are confounded by the
closed-source nature of Tinker's implementation. See
Section~\ref{sec:tinker_closed} for detail.

\paragraph{Cross-platform hardware confounding.}
Tinker, Modal, Colab, and the PES A100 node use different accelerators,
different memory hierarchies, and different inference/training stacks
(Tinker proprietary, Modal vLLM, Colab PyTorch-native, TRL+Accelerate).
Observed differences in reward dynamics are not purely algorithmic.

\paragraph{Carbon estimates are approximate.}
As discussed in Section~\ref{sec:carbon}, Tinker GPU-hour and TDP are
inferred, not measured. Grid intensity is region-average, not marginal.
Numbers in Table~\ref{tab:carbon2} should be treated as
order-of-magnitude.

\paragraph{Exploratory, not confirmatory.}
This is an exploratory case study, not a pre-registered confirmatory
experiment. We did not pre-register primary hypotheses. All $p$-values are
un-corrected for the number of comparisons actually made across the
paper; the Bonferroni-surviving comparisons in
Section~\ref{sec:results} are so flagged, but most of our secondary
analyses are descriptive.

\paragraph{Geographic and demographic limits.}
All authors are based at a single institution in Bengaluru, India. Our
choice of benchmarks, prompt phrasings, and evaluation design reflects
that context. Independent replication by groups in other regions, on
non-English tasks, and with non-Qwen / non-Llama families is an open
question.

% =============================================================================
% End of ethics_statement.tex
